%%
%% Test: Typography Configuration — v1.1
%% Stage: typography refactor (stages 1-4)
%% 
%% Purpose: Verify typography settings:
%%   - Kashida (Persian letter stretching) via xepersian-hm
%%   - Microtype protrusion
%%   - Line breaking settings
%%   - Widow/orphan control
%%   - Custom column types (L, C, R)
%% 

\documentclass{matinbook}

\title{تست تایپوگرافی — نسخه ۱.۱}
\author{تیم MatinBook}
\date{۱۴۰۵}

\begin{document}

\maketitle

\chapter{بررسی تایپوگرافی}
\label{chap:test-typography}

این سند، تنظیمات تایپوگرافی \texttt{mb-typography} را بررسی می‌کند.

%----------------------------------------
\section{کشیدگی (Kashida)}
\label{sec:kashida}

متن فارسی برای هم‌ترازی (justify) از کشیدگی حروف استفاده می‌کند، نه از افزایش فاصله بین کلمات. این متن باید با کشیدگی طبیعی نمایش داده شود.

این یک پاراگراف طولانی‌تر است تا اثر کشیدگی را بهتر نشان دهد. در تایپوگرافی فارسی، حروف به‌جای فاصله‌گذاری بین کلمات، کشیده می‌شوند تا خطوط هم‌تراز شوند. این روش زیبایی بیشتری به متن فارسی می‌دهد.

%----------------------------------------
\section{شکست خط}
\label{sec:line-break}

تنظیمات شکست خط:
\begin{itemize}
    \item \verb|\emergencystretch| = 1em
    \item \verb|\tolerance| = 3000
    \item \verb|\hbadness| = 3000
    \item \verb|\widowpenalty| = 10000
    \item \verb|\clubpenalty| = 10000
\end{itemize}

%----------------------------------------
\section{ستون‌های سفارشی جدول}
\label{sec:custom-columns}

جدول با ستون‌های سفارشی \texttt{L}, \texttt{C}, \texttt{R}:

\begin{table}[h]
\centering
\caption{جدول با ستون‌های سفارشی}
\label{tab:typography-test}
\begin{tabular}{L{4cm} C{2cm} R{3cm}}
\toprule
\textbf{ستون چپ‌چین} & \textbf{ستون وسط‌چین} & \textbf{ستون راست‌چین} \\
\midrule
متن طولانی برای تست چپ‌چین & ۱۰ & راست‌چین \\
متن کوتاه & ۲۰ & ۲۰۰ \\
\bottomrule
\end{tabular}
\end{table}

%----------------------------------------
\section{ریاضیات}
\label{sec:math}

فرمول ریاضی با فاصله‌گذاری مناسب:

\begin{equation}
    \label{eq:typography-test}
    f(x) = \int_{0}^{\infty} e^{-t^2} \, dt = \frac{\sqrt{\pi}}{2}
\end{equation}

%----------------------------------------
\section{پاراگراف طولانی}
\label{sec:long-paragraph}

این یک پاراگراف طولانی است که برای تست مکانیزم کشیدگی و شکست خط در متن فارسی نوشته شده است. همان‌طور که می‌دانید، تایپوگرافی فارسی تفاوت‌های بنیادینی با تایپوگرافی لاتین دارد. در تایپوگرافی لاتین، برای هم‌ترازی خطوط از افزایش فاصله بین کلمات استفاده می‌شود، در حالی که در تایپوگرافی فارسی، حروف داخل کلمات کشیده می‌شوند. این تفاوت باعث می‌شود که متن فارسی طبیعی‌تر و زیباتر به نظر برسد. پکیج \texttt{xepersian-hm} این قابلیت را برای ما فراهم می‌کند.

%----------------------------------------
\section{نتیجه‌گیری}
\label{sec:conclusion}

اگر این سند بدون خطا کامپایل شود:

\begin{itemize}
    \item ✅ کشیدگی (Kashida) فعال است
    \item ✅ تنظیمات microtype درست است
    \item ✅ شکست خط کار می‌کند
    \item ✅ ستون‌های سفارشی \texttt{L}, \texttt{C}, \texttt{R} کار می‌کنند
    \item ✅ ریاضیات با فاصله‌گذاری درست نمایش داده می‌شود
\end{itemize}

\textbf{وضعیت تست:} قبول

\end{document}
